\section{The exact Blaschke correction in the projected ideal trace}
\label{sec:projected-ideal-blaschke}

At polynomial weight zero the actual ideal projector has an
explicit Hardy-space form.  This makes its contribution to the
trace a concrete, absolutely convergent kernel integral.  The
cross-weight nuclear comparison in MA-MO-014 then gives the
same value at every polynomial exponent.  The correction below
is part of the projected ideal trace; it is not silently set
to zero or identified with the full zero-quotient trace.

Keep all hypotheses of MA-MO-012.  This section allows its
general fixed number field and nontrivial primitive finite-order
character.  Use \(s=0\) initially, put
\(\Pi_a=\{z:-a<\operatorname{Re}z<a\}\), and write
\[
 \begin{gathered}
 \varphi_a(z)=\tan\frac{\pi z}{4a},\qquad
 \mathcal B(z)=B(\varphi_a(z)),\\
 Z(B)=\{\varphi_a(\rho-1/2):
                    |\operatorname{Re}(\rho-1/2)|<a\},\\
 b_\pm(t)=\mathcal B(\pm a+it).
 \end{gathered}
 \tag{MA-MO.230}
\]
Multiplicities are retained, \(B\) is the disk Blaschke product
of MA-HY-001, and \(B=1\) when the interior divisor is empty.
Its irrelevant unimodular normalization will cancel.  The
two regimes of MA-MO-012 have no zeros on either edge.

Let \(\mathcal K_b\) denote the full operator
\(V_bT_{\eta,b}(g)V_b^{-1}\), with its Schwartz kernel
\(k_b(t,u)\), for \(b=\pm a\).  This includes all terms of
(MA-MO.187), in particular the compact residual when the
original cutoff is not symmetric.  Define
\[
 l_{\pm a}(t,u)=\overline{b_\pm(t)}\,k_{\pm a}(t,u)b_\pm(u).
 \tag{MA-MO.231}
\]

\paragraph{MA-MO-015 (an explicit projected-trace correction).}
The kernels in (MA-MO.231) are Schwartz.  For every \(s\ge0\),
the actual completed-ideal endomorphism has trace
\[
 \begin{aligned}
 \operatorname{Tr}_{J_s}B_{g,s}
  ={}&\frac12\sum_{\epsilon=\pm1}
              \int_{\mathbb R} k_{\epsilon a}(t,t)\,dt
                         +\mathfrak C_{\eta,a}(g,\phi),\\
 \mathfrak C_{\eta,a}(g,\phi)
  ={}&\frac{i}{8a}\sum_{\epsilon=\pm1}\epsilon
       \int_{\mathbb R}\int_0^\infty
       \frac{l_{\epsilon a}(t,t+v)-l_{\epsilon a}(t,t-v)}
                         {\sinh(\pi v/(4a))}\,dv\,dt .
 \end{aligned}
 \tag{MA-MO.232}
\]
Both displayed integrals are absolutely convergent.  The first
term is one half of the sum of the two unrestricted one-sided
traces.  The second is the exact correction required by the
inherited diagonal-ideal projection.  No vanishing assertion
for that correction is part of this theorem.

\begin{proof}
\emph{The actual strip ideal and its boundary unitary.}
Take the centering parameter of MA-HY to be \(1/2\), its
two strip edges to be \(-a,a\), and its polynomial exponent
to be zero.  Its onto Mellin theorem and its half-derivative
conformal unitary give
\[
 \mathcal M H_0=\mathscr H^2(\Pi_a),\qquad
 \mathcal M J_0=\mathcal B\mathscr H^2(\Pi_a).
 \tag{MA-MO.233}
\]
The notation \(J_0\) in this section means the exponent-zero
ideal, not the inversion operator of the original source.
Equation (MA-MO.233) is the actual closure theorem, not a
new choice of a spectral kernel.  The half-derivative factor
commutes with multiplication as in MA-HY.7, so it introduces
no further factor into \(\mathcal B\).

Boundary values satisfy \(|b_\pm(t)|=1\) almost everywhere.
Thus \(U_{\mathcal B}=\operatorname{diag}(M_{b_+},M_{b_-})\)
is unitary on the two boundary \(L^2\) spaces.  Let
\(\widehat\Pi_D\) be the projection onto the Mellin boundary
image of the entire diagonal \(H_0\), and let
\(\widehat\Pi_J\) project onto its actual ideal image.
Equation (MA-MO.233) implies the equality of orthogonal
projections
\[
          \widehat\Pi_J
            =U_{\mathcal B}\widehat\Pi_D U_{\mathcal B}^*.
 \tag{MA-MO.234}
\]
This retains the whole ideal projector.

\emph{Boundary regularity with the exact generic scope.}
We give the needed derivative bound explicitly.  Write an
interior centered zero as \(\lambda=\beta+i\gamma\), put
\(\kappa=\pi/(2a)\), and map the strip to the upper half-plane
by \(w(z)=\exp(i\kappa(z+a))\).  Up to a constant phase,
the elementary Blaschke factor is
\((w(z)-w(\lambda))/(w(z)-\overline{w(\lambda)})\).
The magnitude of its boundary logarithmic derivative is
\[
 \begin{aligned}
 p_{\lambda,+}(t)
    &=\frac{\kappa\sin(\kappa(\beta+a))}
        {\cosh(\kappa(t-\gamma))+
                              \cos(\kappa(\beta+a))},\\
 p_{\lambda,-}(t)
    &=\frac{\kappa\sin(\kappa(\beta+a))}
        {\cosh(\kappa(t-\gamma))-
                              \cos(\kappa(\beta+a))}.
 \end{aligned}
 \tag{MA-MO.235}
\]
These follow by differentiating the elementary factor on
the two real boundary rays in the \(w\)-plane.  The signs
of the logarithmic derivatives are immaterial for the
following absolute estimates.

For \(|t-\gamma|\le1\), their \(j\)th derivatives are
bounded by \(C_{a,j}\delta_{\lambda,\pm}^{-j-1}\), where
\(\delta_{\lambda,+}=a-\beta\) and
\(\delta_{\lambda,-}=a+\beta\).  Indeed the denominator
is comparable near the relevant edge to
\((t-\gamma)^2+\delta_{\lambda,\pm}^2\); away from that
edge it is bounded below.  For \(|t-\gamma|\ge1\), all
these derivatives are bounded by
\(C_{a,j}e^{-\kappa|t-\gamma|}\), uniformly in \(\beta\).
The local completed-zero count, with multiplicity, is
\(O(\log(2+|T|))\), as in the analytic inputs and
(M18.2) of MT-M18.

In the generic regime, the distance condition and functional
reflection used in MA-MO-012 give
\(\delta_{\lambda,\pm}\ge c(1+|\gamma|)^{-2}\).
In the gap regime they give a fixed positive lower bound.
Summing (MA-MO.235) and its differentiated series therefore
gives polynomial bounds in \(t\) for every derivative of
the boundary logarithmic derivative.  The far-height sums
converge locally uniformly by their exponential tails.
Finite Blaschke products and their derivatives converge
across each finite boundary segment: no zero accumulates
there, and the same logarithmic-derivative estimates apply.
Consequently the asserted boundary values are smooth, and
\[
        |\partial_t^j b_\pm(t)|
           \le C_j(1+|t|)^{N_j}\qquad(j\ge0).
 \tag{MA-MO.236}
\]
One can obtain this last bound recursively from
\(b_\pm'=b_\pm(\log b_\pm)'\) and \(|b_\pm|=1\).
This proves that multiplying the Schwartz kernel \(k_b\)
by the two factors in (MA-MO.231) preserves Schwartz
space.  It does not assume an edge separation at widths
outside the two regimes under consideration.

\emph{The inherited diagonal projection.}
Before Fourier transform, the unweighted two-component
embedding sends \(f\) to \((e^{ax}f,e^{-ax}f)\).
Its pointwise orthogonal projection is
\[
 \Pi_D(x)=\frac1{e^{2ax}+e^{-2ax}}
       \begin{pmatrix}e^{2ax}&1\\1&e^{-2ax}\end{pmatrix},
 \qquad
 \theta_\pm(x)=\frac{e^{\pm2ax}}{e^{2ax}+e^{-2ax}}.
 \tag{MA-MO.237}
\]
Thus the two diagonal blocks of \(\widehat\Pi_D\) are
\(\mathcal F_+M_{\theta_\pm}\mathcal F_+^{-1}\), where
\(\mathcal F_+\) is unitary Fourier transform with the
positive exponential.  Using (MA-MO.190), (MA-MO.234)
and bounded trace cyclicity gives the exact identity
\[
 \operatorname{Tr}_{J_0}B_{g,0}
   =\sum_{\epsilon=\pm1}
       \operatorname{Tr}_{L^2}\!\left(
          \mathcal L_{\epsilon a}\,
          \mathcal F_+M_{\theta_\epsilon}\mathcal F_+^{-1}
                                \right),
 \tag{MA-MO.238}
\]
where \(\mathcal L_b\) has kernel \(l_b\).  Off-diagonal
blocks of \(\widehat\Pi_D\) have not been declared zero;
they simply do not enter the trace of a product whose
other factor is block diagonal.

\emph{Fourier evaluation and the sign.}
For this calculation only write
\(\widehat f_+(v)=\int e^{ivx}f(x)\,dx\), without
the unitary scalar.  Since
\(\theta_+'(x)=a\operatorname{sech}^2(2ax)\), elementary
contour integration, or the beta integral after
\(y=e^{4ax}\), gives
\[
 \begin{aligned}
 \widehat{\theta_+'}_+(v)
       &=\frac{\pi v}{4a\sinh(\pi v/(4a))},\\
 \widehat{\theta_\pm}_+(v)
       &=\pi\delta_0(v)
           \ \pm\frac{i\pi}{4a}\operatorname{PV}
                       \frac1{\sinh(\pi v/(4a))}.
 \end{aligned}
 \tag{MA-MO.239}
\]
Indeed the substitution in the first line gives
\(\int_0^\infty y^{iv/(4a)}(1+y)^{-2}\,dy
=\Gamma(1+iv/(4a))\Gamma(1-iv/(4a))\), and Euler's
reflection identity gives the displayed value, including its
continuous value at zero.  The derivative convention is
\(\widehat{f'}_+=-iv\widehat f_+\); the delta coefficient
is fixed by \(\theta_+(x)+\theta_+(-x)=1\).

For a Schwartz kernel \(l\), the physical kernel obtained
by inverse double Fourier transform is Schwartz.  Hence
the trace against the bounded multiplier \(\theta_\pm\)
is its absolutely integrable diagonal pairing.  Equivalently
it is the distributional pairing
\[
 \operatorname{Tr}\bigl(\mathcal L\mathcal F_+
                    M_{\theta_\pm}\mathcal F_+^{-1}\bigr)
     =\frac1{2\pi}\int\!\!\int
             l(t,u)\widehat{\theta_\pm}_+(u-t)\,dt\,du.
 \tag{MA-MO.240}
\]
The order \(u-t\), rather than \(t-u\), comes from the
second kernel evaluated at \((u,t)\) in the trace product.
Substituting (MA-MO.239), and pairing positive and negative
\(v=u-t\), yields
\[
 \frac12\int l(t,t)\,dt
       \ \pm\frac{i}{8a}\int_{\mathbb R}\int_0^\infty
          \frac{l(t,t+v)-l(t,t-v)}{\sinh(\pi v/(4a))}
                                                  \,dv\,dt.
 \tag{MA-MO.241}
\]
Near zero, the numerator is \(O(v)\), with an integrable
Schwartz majorant in \(t\), by the mean value formula.
For \(v\ge1\), the denominator gives exponential decay,
and the Schwartz bounds supply an integrable majorant.
The second integral is therefore absolutely convergent.
Since \(l_b(t,t)=k_b(t,t)\), (MA-MO.238)--(MA-MO.241)
prove (MA-MO.232) at \(s=0\).  The actual cross-weight
cyclicity (MA-MO.212) proves it at every \(s\ge0\).
\end{proof}

\paragraph{Scope of the explicit correction.}
The kernels \(k_b\) are the explicit original-source
Schwartz kernels of MA-MO-012.  The product \(\mathcal B\)
encodes precisely the retained interior divisor through the
already proved actual closure theorem MA-HY-001.  Thus
(MA-MO.232) gives an explicit projected-trace correction,
not a replacement of the projector by an assigned zero sum.
Its finite theta Gram evaluation and independent matrix
error certificate are MA-MO-014 for the two requested
character families.

The proof uses the exact Mellin and conformal normalization
MA-HY.3--.9, including the independent reproducing-kernel
verification MA-HY.6a--.6b, and the full kernel construction
and genuine ideal range of MA-MO-012.  The local zero count
is the inherited MT-M18 analytic input, not a newly assumed
spacing hypothesis.  The Fourier constant and the diagonal
sum-norm projector were derived explicitly above.

The statement does not identify the result with the
unrestricted sum of the one-sided traces, the original
positive-source trace, or the full zero-quotient trace.
A further genuine nuclear comparison can establish such
an identification if its domains and common trace pairing
are proved; density of the smooth ideal alone cannot.
No arithmetic assertion that an exterior divisor is absent,
and no additional motivic or weight-theoretic positivity,
is inferred from the finite-length divisor or this Hilbert
trace calculation.
